%%%%%%%%%%%%%%%%%%%%%%%%%%%%%%%%%%%%%%%%%%%%%%%%%%%%%%%%%%%%%%%%%%%%
% bare_jrnl.tex
% V1.4b
% 2015/08/26
% by Michael Shell
% see http://www.michaelshell.org/
% for current contact information.                                                                 
%%%%%%%%%%%%%%%%%%%%%%%%%%%%%%%%%%%%%%%%%%%%%%%%%%%%%%%%%%%%%%%%%%%%

%% This is a skeleton file demonstrating the use of IEEEtran.cls
%% (requires IEEEtran.cls version 1.8b or later) with an IEEE
%% journal paper.
%%
%% Support sites:
%% http://www.michaelshell.org/tex/ieeetran/
%% http://www.ctan.org/pkg/ieeetran
%% and
%% http://www.ieee.org/


\documentclass[journal]{IEEEtran}
\raggedbottom
%
\usepackage{instructivo}  % Commands for multiple choice
\graphicspath{{./}{./fig/}}

\usepackage[natbib=true,style=trad-unsrt,backend=biber]{biblatex}

\addbibresource{literatura.bib}
\usepackage{multirow}
\usepackage{graphicx}
\usepackage{array}
\usepackage{float}
\usepackage{capt-of}
\newcolumntype{C}{>{\centering\arraybackslash}X}
% correct bad hyphenation here
\hyphenation{op-tical net-works semi-conduc-tor}


\begin{document}
%
% paper title
% Titles are generally capitalized except for words such as a, an, and, as,
% at, but, by, for, in, nor, of, on, or, the, to and up, which are usually
% not capitalized unless they are the first or last word of the title.
% Linebreaks \\ can be used within to get better formatting as desired.
% Do not put math or special symbols in the title.
\title{Experimento 6: Amplificador de potencia clase B}


\author{Gerald~David~Blanco-Sáenz,~\IEEEmembership{Estudiante,~ITCR}
        y~Melanie~Espinoza-Hernandez,~\IEEEmembership{Estudiante,~ITCR,}}


% The paper headers
\markboth{EL3215 Laboratorio de Electrónica Analógica, IS2026}%
{EL3215 Laboratorio de Electrónica Analógica}


% make the title area
\maketitle


\begin{abstract}
En este experimento se estudió el funcionamiento de un amplificador de potencia clase B en configuración push-pull utilizando transistores bipolares complementarios. Inicialmente se analizaron las formas de onda de entrada y salida del amplificador, observando la presencia de distorsión por cruce cuando los transistores operan cerca del punto de corte. Posteriormente se implementó un circuito de polarización mediante diodos con el objetivo de reducir esta distorsión y mejorar la linealidad de la señal amplificada. Se calcularon y midieron diferentes parámetros de corriente continua y corriente alterna; y los resultados obtenidos permitieron verificar el comportamiento teórico del amplificador clase B y su mayor eficiencia en comparación con amplificadores clase A.
\end{abstract}

% Se sugiere no más de cuatro palabras o frases cortas en orden alfabético, separadas por comas, que representen su reporte
\begin{IEEEkeywords}
Amplificador, potencia, push-pull, clase B.
\end{IEEEkeywords}


%%%%%%%%%%%%%%%%%%%%%%%%%%%%%%%%%%%%%%%%%%%%%%%%%%%%%%%%%%%%%%%%%%%%%%%%%%%%%%%%%%%%%%%%%%%%%%
%%%%%%%%%%%%%%%%%%%%%%%%%%%%%%%%%%%%%%%%%%%%%%%%%%%%%%%%%%%%%%%%%%%%%%%%%%%%%%%%%%%%%%%%%%%%%%
\section{Introducción}

\IEEEPARstart{L} os amplificadores de potencia son circuitos diseñados para entregar una cantidad significativa de energía a una carga, como un parlante o un sistema de audio. A diferencia de los amplificadores de señal pequeña, estos circuitos deben considerar aspectos importantes como la eficiencia energética y la disipación de potencia en los dispositivos activos. Este tipo de amplificadores utilizan una configuración push-pull formada por transistores complementarios que conducen en semiciclos opuestos de la señal, lo que permite mejorar la eficiencia del circuito en comparación con los amplificadores clase A.\cite{Floyd2008} \\
Sin embargo, los amplificadores clase B presentan un fenómeno conocido como distorsión por cruce, que ocurre cuando ninguno de los transistores conduce cerca del punto donde la señal cruza por cero debido a la tensión necesaria para polarizar la unión base-emisor. En este experimento se construyó y analizó un amplificador de potencia clase B en configuración push-pull, observando las formas de onda de entrada y salida, la presencia de distorsión por cruce y el efecto de utilizar diodos de polarización para reducir dicha distorsión y mejorar el comportamiento del amplificador.

%%%%%%%%%%%%%%%%%%%%%%%%%%%%%%%%%%%%%%%%%%%%%%%%%%%%%%%%%%%%%%%%%%%%%%%%%%%%%%%%%%%%%%%%%%%%%%
%%%%%%%%%%%%%%%%%%%%%%%%%%%%%%%%%%%%%%%%%%%%%%%%%%%%%%%%%%%%%%%%%%%%%%%%%%%%%%%%%%%%%%%%%%%%%%
\section{Amplificador de potencia clase B}

\subsection{Circuitos de Medición}
El circuito utilizado para analizar el amplificador clase B se muestra en la Fig.\ref{fig:circuito1}. Este corresponde a una configuración push-pull con transistores complementarios, donde cada transistor amplifica un semiciclo de la señal de entrada y la salida se obtiene sobre la resistencia de carga RL.

\begin{figure}[H]
	\centering
	\includegraphics[width=3in]{fig/Circuito clase B_Version 1.jpg}
	\caption{Amplificador de potencia clase B}
	\label{fig:circuito1}
\end{figure}


Posteriormente se implementó una segunda versión del circuito como el de la Fig.\ref{fig:circuito2} que incorpora diodos de polarización, los cuales generan una caída de tensión aproximada de 0,7 V en cada unión base-emisor. Esto permite que ambos transistores se encuentren ligeramente polarizados, reduciendo el efecto de distorsión por cruce presente en el amplificador clase B ideal.

\begin{figure}[H]
	\centering
	\includegraphics[width=2.8in]{fig/Circuito clase B_Version 2.jpg}
	\caption{Amplificador de potencia clase B: Versión 2}
	\label{fig:circuito2}
\end{figure}

Finalmente se analizó una tercera configuración del amplificador como se presenta en la Fig.\ref{fig:circuito3} , en la cual se agregó un transistor adicional para proporcionar ganancia de tensión al circuito.

\begin{figure}[H]
	\centering
	\includegraphics[width=3in]{fig/Circuito clase B_Version 3.jpg}
	\caption{Amplificador de potencia clase B: Versión 3}
	\label{fig:circuito3}
\end{figure}


\subsection{Simulaciones}
Inicialmente se simuló el circuito mostrado en la Fig.\ref{fig:circuito1}, observando las formas de onda de entrada y salida. 

\begin{figure}[H]
	\centering
	\includegraphics[width=3in]{fig/Onda de entrada y salida versión 1.png}
	\caption{Onda entrada (roja) y salida (azul) del circuito de la Figura \ref{fig:circuito1} }
	\label{fig:ondas1}
\end{figure}

Como se aprecia en la Fig.\ref{fig:ondas1}, la señal de salida sigue la forma de la señal de entrada pero presenta una pequeña distorsión cerca del cruce por cero, característica de los amplificadores clase B debido a la tensión mínima necesaria para polarizar las uniones base-emisor de los transistores.\\

Posteriormente se simuló el circuito mostrado en la Fig.\ref{fig:circuito2}, a partir de esta configuración se calcularon los parámetros de corriente continua mostrados en la Tabla.\ref{tab:cd_version2} y corriente alterna mostrados en la Tabla.\ref{tab:ca_version2} del circuito.

\begin{table}[H]
\centering
\caption{Amplificador de potencia clase B versión 2: Parámetros CD}
\label{tab:cd_version2}
\begin{tabular}{|c|c|}
\hline
Parámetros CD & Valor calculado  \\
\hline
$V_{E}$ & $0\ V$ \\
\hline
$V_{B1}$ & $-0.745\ V$ \\
\hline
$V_{B2}$ & $0.931\ V$ \\
\hline
$I_{R1} = I_{CQ}$ & $825\ \mu A$ \\
\hline
\end{tabular}
\end{table}

\begin{table}[H]
\centering
\caption{Amplificador de potencia clase B versión 2: Parámetros CA}
\label{tab:ca_version2}
\begin{tabular}{|c|c|}
\hline
Parámetros CA & Valor calculado  \\
\hline
$V_{p (sal)}$ & $3.5\ V$ \\
\hline
$V_{p (sal)}$ & $10.6\ mA$ \\
\hline
$P_{sal}$ & $18.56\ mW$ \\
\hline
\end{tabular}
\end{table}

Finalmente se analizó el circuito de la Fig.\ref{fig:circuito3} , en el cual se agrega un transistor adicional que funciona como etapa de amplificación previa. En esta configuración se calcularon los parámetros de polarización del transistor Q3 mostrados en la Tabla.\ref{tab:cd_version3} y la ganancia de tensión aproximada del amplificador completo en la Tabla.\ref{tab:ca_version3}.

\begin{table}[H]
\centering
\caption{Amplificador de potencia clase B versión 3: Parámetros CD}
\label{tab:cd_version3}
\begin{tabular}{|c|c|}
\hline
Parámetros CD & Valor calculado  \\
\hline
$V_{B (Q3)}$ & $-6.36\ V$ \\
\hline
$V_{E (Q3)}$ & $-7.01\ V$ \\
\hline
$I_{C (Q3)}$ & $731\ \mu A$ \\
\hline
\end{tabular}
\end{table}

\begin{table}[H]
\centering
\caption{Amplificador de potencia clase B versión 3: Parámetros CA}
\label{tab:ca_version3}
\begin{tabular}{|c|c|}
\hline
Parámetros CA & Valor calculado  \\
\hline
$A'_{V (Q3)}$ & $2.0745$ \\
\hline
\end{tabular}
\end{table}


\subsection{Resultados Experimentales}

Una vez realizadas las simulaciones, se procedió a construir físicamente el circuito del amplificador clase B y a registrar los valores experimentales correspondientes. \\

En primer lugar se midieron las resistencias utilizadas para verificar que sus valores se encontraran dentro de las tolerancias esperadas.

\begin{table}[H]
\centering
\caption{Amplificador de potencia clase B versión 2: Parámetros CA}
\label{tab:resistencias}
\begin{tabular}{|c|c|c|}
\hline
Componente & Valor requerido & Valor Medido  \\
\hline
$R_{L}$ & $330\ \Omega $ & $  330.1 \Omega $ \\
\hline
$R_{1}$ & $10\ k\Omega $ & $  9.995 k\Omega $ \\
\hline
$R_{2}$ & $10\ k\Omega $ & $  10.001 k\Omega $ \\
\hline
$R_{3}$ & $68\ k\Omega $ & $  67.8 k\Omega $ \\
\hline
$R_{4}$ & $2.7\ k\Omega $ & $  2.7 k\Omega $ \\
\hline
\end{tabular}
\end{table}

Luego, se evaluó el circuito de la Fig.\ref{fig:circuito1}. En esta configuración se observaron las señales de entrada y salida utilizando el osciloscopio para verificar el comportamiento real del amplificador, aunque se presentó una pequeña distorsión cerca del cruce por cero.

\begin{figure}[H]
	\centering
	\includegraphics[width=3in]{fig/Ondas experimentales.jpeg}
	\caption{Ondas entrada y salida experimentales del circuito de la Figura \ref{fig:circuito1} }
	\label{fig:ondas2}
\end{figure}

Posteriormente se implementó el circuito de la Fig.\ref{fig:circuito2} que incorpora los diodos. Con esta modificación se realizaron las mediciones de los parámetros de corriente continua y corriente alterna del amplificador.

\begin{table}[H]
\centering
\caption{Amplificador de potencia clase B versión 2: Parámetros CD experimentales}
\label{tab:cd_version2_exp}
\begin{tabular}{|c|c|}
\hline
Parámetros CD & Valor medido  \\
\hline
$V_{E}$ & $0.001\ V$ \\
\hline
$V_{B1}$ & $0.5967\ V$ \\
\hline
$V_{B2}$ & $0.5969\ V$ \\
\hline
$I_{R1} = I_{CQ}$ & \cellcolor{gray!85} \\
\hline
\end{tabular}
\end{table}

\begin{table}[H]
\centering
\caption{Amplificador de potencia clase B versión 2: Parámetros CA experimentales}
\label{tab:ca_version2_exp}
\begin{tabular}{|c|c|}
\hline
Parámetros CA & Valor medido  \\
\hline
$V_{p (sal)}$ & $7.3\ V$ \\
\hline
$V_{p (sal)}$ & \cellcolor{gray!85} \\
\hline
$P_{sal}$ & \cellcolor{gray!85} \\
\hline
\end{tabular}
\end{table}

Finalmente se analizó el circuito de la Fig.\ref{fig:circuito3}, en el cual se agrega un transistor adicional como etapa de amplificación previa. En esta configuración se midieron los parámetros de polarización del transistor Q3 y la ganancia de tensión aproximada del amplificador.

\begin{table}[H]
\centering
\caption{Amplificador de potencia clase B versión 3: Parámetros CD experimentales}
\label{tab:cd_version3_exp}
\begin{tabular}{|c|c|}
\hline
Parámetros CD & Valor calculado  \\
\hline
$V_{B (Q3)}$ & $-6.197\ V$ \\
\hline
$V_{E (Q3)}$ & $-7.157\ V$ \\
\hline
$I_{C (Q3)}$ & \cellcolor{gray!85} \\
\hline
\end{tabular}
\end{table}

\begin{table}[H]
\centering
\caption{Amplificador de potencia clase B versión 3: Parámetros CA experimentales}
\label{tab:ca_version3_exp}
\begin{tabular}{|c|c|}
\hline
Parámetros CA & Valor calculado  \\
\hline
$A'_{V (Q3)}$ & $1.15$ \\
\hline
\end{tabular}
\end{table}


\subsection{Análisis de Resultados}
En este tipo de amplificadores cada transistor conduce aproximadamente durante $180^\circ$ del ciclo de la señal, lo que permite mejorar la eficiencia respecto a los amplificadores clase A, aunque introduce el fenómeno conocido como distorsión por cruce. \cite{Floyd2008}\\

En la primera configuración del circuito (Fig.\ref{fig:circuito1}) se implementó un amplificador clase B básico utilizando transistores complementarios. En esta configuración el transistor NPN conduce durante el semiciclo positivo de la señal de entrada mientras que el transistor PNP conduce durante el semiciclo negativo. Sin embargo, debido a que cada transistor requiere aproximadamente $0.7\,V$ para polarizar su unión base-emisor, existe una región cercana al cruce por cero en la que ninguno de los transistores conduce, produciendo una distorsión en la señal de salida conocida como \textit{distorsión por cruce}.\cite{Boylestad} \\

Para reducir este efecto se implementó la segunda configuración del circuito (Fig.\ref{fig:circuito2}), en la cual se añadieron dos diodos de polarización entre las bases de los transistores. Estos diodos generan una caída de tensión aproximada de $1.4\,V$, equivalente a la suma de las dos caídas base-emisor, lo que permite que los transistores permanezcan ligeramente polarizados en conducción incluso cuando la señal de entrada se encuentra cerca de cero. Como resultado, la transición entre semiciclos se vuelve más continua y la distorsión por cruce se reduce considerablemente.\\

A partir de esta configuración se determinaron los parámetros de operación del circuito. El voltaje de emisor se mantiene cercano a $0\,V$, lo cual es consistente con el uso de fuentes simétricas de $\pm 9\,V$, permitiendo que la señal de salida oscile alrededor de cero.\cite{SedraSmith2015} Asimismo, los voltajes de base obtenidos se aproximan a $V_{B1}=-0.745\,V$ y $V_{B2}=0.931\,V$, valores que corresponden a las caídas de tensión generadas por los diodos de polarización.

En cuanto al comportamiento en señal, se obtuvo un voltaje pico de salida aproximado de

\[
V_{p(sal)} = 3.5\,V
\]

lo que produce una corriente pico en la carga dada por

\[
I_{p(sal)} = \frac{V_p}{R_L}
\]

considerando $R_L = 330\,\Omega$:

\[
I_{p(sal)} = \frac{3.5}{330} \approx 10.6\,mA
\]

Este resultado de $V_{p(sal)}$ coincide con el valor obtenido experimentalmente. A partir de estos datos también se puede estimar la potencia entregada a la carga mediante

\[
P_{sal} = \frac{V_{rms}^2}{R_L}
\]

donde

\[
V_{rms} = \frac{V_p}{\sqrt{2}}
\]

lo cual conduce a una potencia aproximada de $18.56\,mW$, valor consistente con las mediciones realizadas durante el experimento. \\

Finalmente se analizó la tercera configuración del amplificador (Fig.\ref{fig:circuito3}), en la cual se incorporó el transistor $Q_3$ como una etapa de amplificación previa. A partir del análisis de polarización se obtuvo un voltaje de base de aproximadamente $-6.36\,V$ y un voltaje de emisor de $-7.01\,V$, con una corriente de colector cercana a $731\,\mu A$, lo cual indica que el transistor opera en la región activa.

La ganancia de tensión aproximada de esta etapa se determinó como

\[
A'_v \approx 2.07
\]

lo cual confirma que el transistor adicional proporciona una amplificación moderada de la señal antes de la etapa de potencia. No obstante, la ganancia total del circuito sigue siendo relativamente baja debido a que la etapa  \textit{push-pull} funciona principalmente como etapa de potencia, cuyo objetivo principal no es aumentar significativamente el voltaje sino entregar mayor corriente a la carga.\cite{Floyd2008}\\

En general, los resultados experimentales obtenidos muestran un comportamiento coherente con el funcionamiento teórico de los amplificadores clase B descritos en la literatura. La incorporación de diodos de polarización permitió reducir la distorsión por cruce, mientras que la inclusión de una etapa previa de amplificación permitió aumentar la amplitud de la señal aplicada a la etapa de salida.
%%%%%%%%%%%%%%%%%%%%%%%%%%%%%%%%%%%%%%%%%%%%%%%%%%%%%%%%%%%%%%%%%%%%%%%%%%%%%%%%%%%%
%%%%%%%%%%%%%%%%%%%%%%%%%%%%%%%%%%%%%%%%%%%%%%%%%%%%%%%%%%%%%%%%%%%%%%%%%%%%%%%%%%%%


%%%%%%%%%%%%%%%%%%%%%%%%%%%%%%%%%%%%%%%%%%%%%%%%%%%%%%%%%%%%%%%%%%%%%%%%%%%%%%%%%%%%
%%%%%%%%%%%%%%%%%%%%%%%%%%%%%%%%%%%%%%%%%%%%%%%%%%%%%%%%%%%%%%%%%%%%%%%%%%%%%%%%%%%%
\section{Conclusiones}

\begin{enumerate}
    \item  Los resultados obtenidos para la versión polarizada del amplificador clase B fueron coherentes con el comportamiento esperado.

    \item La polarización con diodos mejoró el funcionamiento del amplificador, al establecer tensiones de base cercanas a las caídas esperadas en las uniones y reducir la distorsión por cruce observada en la configuración básica del push-pull.

    \item La incorporación de la etapa con $Q_3$ permitió obtener una ganancia de tensión adicional, evidenciando que esta etapa funciona como preamplificador de la señal. De esta manera, la etapa push-pull queda principalmente encargada de suministrar la corriente y la potencia necesarias a la carga.

\end{enumerate}

%%%%%%%%%%%%%%%%%%%%%%%%%%%%%%%%%%%%%%%%%%%%%%%%%%%%%%%%%%%%%%%%%%%%%%%%%%%%%%%%%%%%
%%%%%%%%%%%%%%%%%%%%%%%%%%%%%%%%%%%%%%%%%%%%%%%%%%%%%%%%%%%%%%%%%%%%%%%%%%%%%%%%%%%%
\section{Referencias Bibliográficas}
\printbibliography


%%%%%%%%%%%%%%%%%%%%%%%%%%%%%%%%%%%%%%%%%%%%%%%%%%%%%%%%%%%%%%%%%%%%%%%%%%%%%%%%%%%%
%%%%%%%%%%%%%%%%%%%%%%%%%%%%%%%%%%%%%%%%%%%%%%%%%%%%%%%%%%%%%%%%%%%%%%%%%%%%%%%%%%%%
\begin{IEEEbiography}[{\includegraphics[width=1in,height=1.25in,clip,keepaspectratio]{Gerald.PNG}}]{Gerald Blanco Sáenz}
Estudiante de Ingeniería en Electrónica desde 2022, graduado del Liceo UNESCO y nacido en Pérez Zeledón. 
Correo: geblanco@estudiantec.cr
\end{IEEEbiography}

\begin{IEEEbiography}[{\includegraphics[width=1in,height=1.25in,clip,keepaspectratio]{Melanie.png}}]{Melanie Espinoza Hernández}
Nació el 1 de junio de 2003 en San José, Costa Rica. Realizó sus estudios en el Colegio Técnico Profesional de Puriscal. 
Se graduó con bachillerato en educación media y obtuvo el título de técnico medio en Computer Science in Software Development. Ingresó al Instituto Tecnológico de Costa Rica y actualmente cursa la carrera de Ingeniería en Electrónica. 
Correo: melespinoza@estudiantec.cr
\end{IEEEbiography}



\vfill

\end{document}






